% --- Archon physics formalization source begin ---
% archon:physics
% archon:covers PhyXMiniProblems/problem_phyx_mini_0893.lean
% archon:source-report reports/phyx_mini/problem_phyx_mini_0893.source.json

\chapter{Physics problem phyx\_mini\_0893}
\label{ch:PhyXMiniProblems_problem_phyx_mini_0893}

\paragraph{Problem source.}
\#\# Question

What is the magnitude of the force \textbackslash{}( \textbackslash{}vec\{F\} \textbackslash{}) on the \textbackslash{}( 1.0\textbackslash{},\textbackslash{}mathrm\{nC\} \textbackslash{}) charge in figure?

\#\# Answer choices

- A: A: 1.25 \textbackslash{}times 10\textasciicircum{}\{-3\}N

- B: B: 1.35 \textbackslash{}times 10\textasciicircum{}\{-3\}N

- C: C: 3.1 \textbackslash{}times 10\textasciicircum{}\{-4\}N

- D: D: 1.33 \textbackslash{}times 10\textasciicircum{}\{-3\}N

\#\# Figure caption (auxiliary; use the image as primary evidence)

The image shows an equilateral triangle with three charges at each vertex. The vertices are labeled with their respective charges:

1. Top vertex: +1.0 nC
2. Bottom left vertex: +2.0 nC
3. Bottom right vertex: +2.0 nC

Each side of the triangle measures 1.0 cm. The internal angles at each vertex are 60 degrees. Dashed lines represent the sides of the triangle, and there is a dotted circle around each charge indicating their positive nature. The distances between charges are annotated alongside the triangle.

\#\# Dataset metadata

Electrostatics; Physical Model Grounding Reasoning, Multi-Formula Reasoning

\paragraph{Recorded answer/context.}
C: C: 3.1 \textbackslash{}times 10\textasciicircum{}\{-4\}N

\paragraph{Figure/image path.}
/root/proposal\_for\_physic/science-mango/phyx\_mini\_run/phyx\_data/test\_image/893.png

\paragraph{Formalization target.}
create a compiling Lean file with sorry bodies at `PhyXMiniProblems/problem\_phyx\_mini\_0893.lean`. The Lean declarations must preserve the physical quantities, dimensions or dimensional roles, figure labels, governing-law hypotheses, and final relation expressed by this problem.
Use Mathlib/Physlib names found through LeanExplore where available. If a physics API is missing, introduce faithful local abstractions rather than scalar placeholder aliases.

\begin{theorem}[Physics formalization target]
\label{thm:physics:phyx_mini_0893:target}
\lean{PhyXMiniProblems.ProblemPhyXMini0893.problem_phyx_mini_0893}
\uses{def:physics:phyx-mini-0893:phyxminiproblems-problemphyxmini0893-equilateralthreechargesetup, def:physics:phyx-mini-0893:phyxminiproblems-problemphyxmini0893-matchesthreepositivepointchargescenario, def:physics:phyx-mini-0893:phyxminiproblems-problemphyxmini0893-matchessuppliedthreechargetrianglefigure, def:physics:phyx-mini-0893:phyxminiproblems-problemphyxmini0893-usesequilateraltrianglegeometry, def:physics:phyx-mini-0893:phyxminiproblems-problemphyxmini0893-hasphysicalthreechargeparameters, def:physics:phyx-mini-0893:phyxminiproblems-problemphyxmini0893-usestextbookvacuumcoulombconstant, def:physics:phyx-mini-0893:phyxminiproblems-problemphyxmini0893-satisfiescoulombforcelawandsuperposition, def:physics:phyx-mini-0893:phyxminiproblems-problemphyxmini0893-netforcemagnitudeinnewtons, def:physics:phyx-mini-0893:phyxminiproblems-problemphyxmini0893-recordeddatasetanswer, def:physics:phyx-mini-0893:phyxminiproblems-problemphyxmini0893-roundstonearestresolution, def:physics:phyx-mini-0893:phyxminiproblems-problemphyxmini0893-isuniqueclosestdisplayedanswer}
The assigned autoformalize agent should translate this physics problem into Lean declarations in the covered file, with theorem and lemma proof bodies written as `by sorry`.
\end{theorem}
\begin{proof}
This is an autoformalization task, not a proof task. Produce statements that can later be proved by the physics prover without weakening the physical model.
\end{proof}

% --- Archon declaration topology begin ---
\section{Declaration topology}
\begin{definition}[force Dimension]
  \label{def:physics:phyx-mini-0893:phyxminiproblems-problemphyxmini0893-forcedimension}
  \lean{PhyXMiniProblems.ProblemPhyXMini0893.forceDimension}
  The physical dimension M L T⁻² of force
\end{definition}
\begin{proof}
  This is the declared model definition of force Dimension.
\end{proof}
\begin{definition}[coulomb Constant Dimension]
  \label{def:physics:phyx-mini-0893:phyxminiproblems-problemphyxmini0893-coulombconstantdimension}
  \lean{PhyXMiniProblems.ProblemPhyXMini0893.coulombConstantDimension}
  \uses{def:physics:phyx-mini-0893:phyxminiproblems-problemphyxmini0893-forcedimension}
  The physical dimension M L³ T⁻² C⁻² of Coulomb's constant
\end{definition}
\begin{proof}
  This is the declared model definition of coulomb Constant Dimension.
\end{proof}
\begin{definition}[Planar Vector]
  \label{def:physics:phyx-mini-0893:phyxminiproblems-problemphyxmini0893-planarvector}
  \lean{PhyXMiniProblems.ProblemPhyXMini0893.PlanarVector}
  A two-dimensional real vector in the plane of the figure
\end{definition}
\begin{proof}
  This is the declared model definition of Planar Vector.
\end{proof}
\begin{definition}[Length Quantity]
  \label{def:physics:phyx-mini-0893:phyxminiproblems-problemphyxmini0893-lengthquantity}
  \lean{PhyXMiniProblems.ProblemPhyXMini0893.LengthQuantity}
  A nonnegative, unit-independent physical length
\end{definition}
\begin{proof}
  This is the declared model definition of Length Quantity.
\end{proof}
\begin{definition}[Charge Magnitude Quantity]
  \label{def:physics:phyx-mini-0893:phyxminiproblems-problemphyxmini0893-chargemagnitudequantity}
  \lean{PhyXMiniProblems.ProblemPhyXMini0893.ChargeMagnitudeQuantity}
  A nonnegative, unit-independent electric-charge magnitude
\end{definition}
\begin{proof}
  This is the declared model definition of Charge Magnitude Quantity.
\end{proof}
\begin{definition}[Force Magnitude Quantity]
  \label{def:physics:phyx-mini-0893:phyxminiproblems-problemphyxmini0893-forcemagnitudequantity}
  \lean{PhyXMiniProblems.ProblemPhyXMini0893.ForceMagnitudeQuantity}
  \uses{def:physics:phyx-mini-0893:phyxminiproblems-problemphyxmini0893-forcedimension}
  A nonnegative, unit-independent magnitude of force
\end{definition}
\begin{proof}
  This is the declared model definition of Force Magnitude Quantity.
\end{proof}
\begin{definition}[Force Vector Quantity]
  \label{def:physics:phyx-mini-0893:phyxminiproblems-problemphyxmini0893-forcevectorquantity}
  \lean{PhyXMiniProblems.ProblemPhyXMini0893.ForceVectorQuantity}
  \uses{def:physics:phyx-mini-0893:phyxminiproblems-problemphyxmini0893-forcedimension, def:physics:phyx-mini-0893:phyxminiproblems-problemphyxmini0893-planarvector}
  A unit-independent planar force vector
\end{definition}
\begin{proof}
  This is the declared model definition of Force Vector Quantity.
\end{proof}
\begin{definition}[Coulomb Constant Quantity]
  \label{def:physics:phyx-mini-0893:phyxminiproblems-problemphyxmini0893-coulombconstantquantity}
  \lean{PhyXMiniProblems.ProblemPhyXMini0893.CoulombConstantQuantity}
  \uses{def:physics:phyx-mini-0893:phyxminiproblems-problemphyxmini0893-coulombconstantdimension}
  A nonnegative, dimensionful value of Coulomb's constant
\end{definition}
\begin{proof}
  This is the declared model definition of Coulomb Constant Quantity.
\end{proof}
\begin{definition}[length Readout]
  \label{def:physics:phyx-mini-0893:phyxminiproblems-problemphyxmini0893-lengthreadout}
  \lean{PhyXMiniProblems.ProblemPhyXMini0893.lengthReadout}
  \uses{def:physics:phyx-mini-0893:phyxminiproblems-problemphyxmini0893-lengthquantity}
  Read a physical length in the coherent length unit selected by units
\end{definition}
\begin{proof}
  This is the declared model definition of length Readout.
\end{proof}
\begin{definition}[charge Magnitude Readout]
  \label{def:physics:phyx-mini-0893:phyxminiproblems-problemphyxmini0893-chargemagnitudereadout}
  \lean{PhyXMiniProblems.ProblemPhyXMini0893.chargeMagnitudeReadout}
  \uses{def:physics:phyx-mini-0893:phyxminiproblems-problemphyxmini0893-chargemagnitudequantity}
  Read a charge magnitude in the coherent charge unit selected by units
\end{definition}
\begin{proof}
  This is the declared model definition of charge Magnitude Readout.
\end{proof}
\begin{definition}[force Magnitude Readout]
  \label{def:physics:phyx-mini-0893:phyxminiproblems-problemphyxmini0893-forcemagnitudereadout}
  \lean{PhyXMiniProblems.ProblemPhyXMini0893.forceMagnitudeReadout}
  \uses{def:physics:phyx-mini-0893:phyxminiproblems-problemphyxmini0893-forcemagnitudequantity}
  Read a force magnitude in the coherent force unit selected by units
\end{definition}
\begin{proof}
  This is the declared model definition of force Magnitude Readout.
\end{proof}
\begin{definition}[force Vector Readout]
  \label{def:physics:phyx-mini-0893:phyxminiproblems-problemphyxmini0893-forcevectorreadout}
  \lean{PhyXMiniProblems.ProblemPhyXMini0893.forceVectorReadout}
  \uses{def:physics:phyx-mini-0893:phyxminiproblems-problemphyxmini0893-planarvector, def:physics:phyx-mini-0893:phyxminiproblems-problemphyxmini0893-forcevectorquantity}
  Read a planar force vector in the coherent force unit selected by units
\end{definition}
\begin{proof}
  This is the declared model definition of force Vector Readout.
\end{proof}
\begin{definition}[coulomb Constant Readout]
  \label{def:physics:phyx-mini-0893:phyxminiproblems-problemphyxmini0893-coulombconstantreadout}
  \lean{PhyXMiniProblems.ProblemPhyXMini0893.coulombConstantReadout}
  \uses{def:physics:phyx-mini-0893:phyxminiproblems-problemphyxmini0893-coulombconstantquantity}
  Read Coulomb's constant in the coherent unit system selected by units
\end{definition}
\begin{proof}
  This is the declared model definition of coulomb Constant Readout.
\end{proof}
\begin{definition}[length In Meters]
  \label{def:physics:phyx-mini-0893:phyxminiproblems-problemphyxmini0893-lengthinmeters}
  \lean{PhyXMiniProblems.ProblemPhyXMini0893.lengthInMeters}
  \uses{def:physics:phyx-mini-0893:phyxminiproblems-problemphyxmini0893-lengthquantity, def:physics:phyx-mini-0893:phyxminiproblems-problemphyxmini0893-lengthreadout}
  Read a physical length in coherent-SI metres
\end{definition}
\begin{proof}
  This is the declared model definition of length In Meters.
\end{proof}
\begin{definition}[length In Centimeters]
  \label{def:physics:phyx-mini-0893:phyxminiproblems-problemphyxmini0893-lengthincentimeters}
  \lean{PhyXMiniProblems.ProblemPhyXMini0893.lengthInCentimeters}
  \uses{def:physics:phyx-mini-0893:phyxminiproblems-problemphyxmini0893-lengthquantity, def:physics:phyx-mini-0893:phyxminiproblems-problemphyxmini0893-lengthinmeters}
  Read a physical length in centimetres
\end{definition}
\begin{proof}
  This is the declared model definition of length In Centimeters.
\end{proof}
\begin{definition}[charge Magnitude In Coulombs]
  \label{def:physics:phyx-mini-0893:phyxminiproblems-problemphyxmini0893-chargemagnitudeincoulombs}
  \lean{PhyXMiniProblems.ProblemPhyXMini0893.chargeMagnitudeInCoulombs}
  \uses{def:physics:phyx-mini-0893:phyxminiproblems-problemphyxmini0893-chargemagnitudequantity, def:physics:phyx-mini-0893:phyxminiproblems-problemphyxmini0893-chargemagnitudereadout}
  Read a charge magnitude in coherent-SI coulombs
\end{definition}
\begin{proof}
  This is the declared model definition of charge Magnitude In Coulombs.
\end{proof}
\begin{definition}[charge Magnitude In Nanocoulombs]
  \label{def:physics:phyx-mini-0893:phyxminiproblems-problemphyxmini0893-chargemagnitudeinnanocoulombs}
  \lean{PhyXMiniProblems.ProblemPhyXMini0893.chargeMagnitudeInNanocoulombs}
  \uses{def:physics:phyx-mini-0893:phyxminiproblems-problemphyxmini0893-chargemagnitudequantity, def:physics:phyx-mini-0893:phyxminiproblems-problemphyxmini0893-chargemagnitudeincoulombs}
  Read a charge magnitude in nanocoulombs
\end{definition}
\begin{proof}
  This is the declared model definition of charge Magnitude In Nanocoulombs.
\end{proof}
\begin{definition}[force Magnitude In Newtons]
  \label{def:physics:phyx-mini-0893:phyxminiproblems-problemphyxmini0893-forcemagnitudeinnewtons}
  \lean{PhyXMiniProblems.ProblemPhyXMini0893.forceMagnitudeInNewtons}
  \uses{def:physics:phyx-mini-0893:phyxminiproblems-problemphyxmini0893-forcemagnitudequantity, def:physics:phyx-mini-0893:phyxminiproblems-problemphyxmini0893-forcemagnitudereadout}
  Read a force magnitude in coherent-SI newtons
\end{definition}
\begin{proof}
  This is the declared model definition of force Magnitude In Newtons.
\end{proof}
\begin{definition}[force Vector In Newtons]
  \label{def:physics:phyx-mini-0893:phyxminiproblems-problemphyxmini0893-forcevectorinnewtons}
  \lean{PhyXMiniProblems.ProblemPhyXMini0893.forceVectorInNewtons}
  \uses{def:physics:phyx-mini-0893:phyxminiproblems-problemphyxmini0893-planarvector, def:physics:phyx-mini-0893:phyxminiproblems-problemphyxmini0893-forcevectorquantity, def:physics:phyx-mini-0893:phyxminiproblems-problemphyxmini0893-forcevectorreadout}
  Read a planar force vector in coherent-SI newtons
\end{definition}
\begin{proof}
  This is the declared model definition of force Vector In Newtons.
\end{proof}
\begin{definition}[coulomb Constant In Newton Meters Squared Per Coulomb Squared]
  \label{def:physics:phyx-mini-0893:phyxminiproblems-problemphyxmini0893-coulombconstantinnewtonmeterssquaredpercoulombsquared}
  \lean{PhyXMiniProblems.ProblemPhyXMini0893.coulombConstantInNewtonMetersSquaredPerCoulombSquared}
  \uses{def:physics:phyx-mini-0893:phyxminiproblems-problemphyxmini0893-coulombconstantquantity, def:physics:phyx-mini-0893:phyxminiproblems-problemphyxmini0893-coulombconstantreadout}
  Read Coulomb's constant in coherent SI N m² / C²
\end{definition}
\begin{proof}
  This is the declared model definition of coulomb Constant In Newton Meters Squared Per Coulomb Squared.
\end{proof}
\begin{definition}[Triangle Vertex]
  \label{def:physics:phyx-mini-0893:phyxminiproblems-problemphyxmini0893-trianglevertex}
  \lean{PhyXMiniProblems.ProblemPhyXMini0893.TriangleVertex}
  The three labelled vertices of the triangular charge configuration
\end{definition}
\begin{proof}
  This is the declared model definition of Triangle Vertex.
\end{proof}
\begin{definition}[Triangle Edge]
  \label{def:physics:phyx-mini-0893:phyxminiproblems-problemphyxmini0893-triangleedge}
  \lean{PhyXMiniProblems.ProblemPhyXMini0893.TriangleEdge}
  The three sides of the triangle, named by their placement in the image
\end{definition}
\begin{proof}
  This is the declared model definition of Triangle Edge.
\end{proof}
\begin{definition}[Lower Charge]
  \label{def:physics:phyx-mini-0893:phyxminiproblems-problemphyxmini0893-lowercharge}
  \lean{PhyXMiniProblems.ProblemPhyXMini0893.LowerCharge}
  The two lower charges that exert force on the requested top charge
\end{definition}
\begin{proof}
  This is the declared model definition of Lower Charge.
\end{proof}
\begin{definition}[vertex]
  \label{def:physics:phyx-mini-0893:phyxminiproblems-problemphyxmini0893-lowercharge-vertex}
  \lean{PhyXMiniProblems.ProblemPhyXMini0893.LowerCharge.vertex}
  \uses{def:physics:phyx-mini-0893:phyxminiproblems-problemphyxmini0893-trianglevertex, def:physics:phyx-mini-0893:phyxminiproblems-problemphyxmini0893-lowercharge}
  Regard a lower-charge label as its corresponding triangle vertex
\end{definition}
\begin{proof}
  This is the declared model definition of vertex.
\end{proof}
\begin{definition}[edge To Top]
  \label{def:physics:phyx-mini-0893:phyxminiproblems-problemphyxmini0893-lowercharge-edgetotop}
  \lean{PhyXMiniProblems.ProblemPhyXMini0893.LowerCharge.edgeToTop}
  \uses{def:physics:phyx-mini-0893:phyxminiproblems-problemphyxmini0893-triangleedge, def:physics:phyx-mini-0893:phyxminiproblems-problemphyxmini0893-lowercharge}
  The side separating a lower source charge from the top charge
\end{definition}
\begin{proof}
  This is the declared model definition of edge To Top.
\end{proof}
\begin{definition}[Charge Sign]
  \label{def:physics:phyx-mini-0893:phyxminiproblems-problemphyxmini0893-chargesign}
  \lean{PhyXMiniProblems.ProblemPhyXMini0893.ChargeSign}
  Sign of an electric charge in the physical model
\end{definition}
\begin{proof}
  This is the declared model definition of Charge Sign.
\end{proof}
\begin{definition}[Electrostatic Environment]
  \label{def:physics:phyx-mini-0893:phyxminiproblems-problemphyxmini0893-electrostaticenvironment}
  \lean{PhyXMiniProblems.ProblemPhyXMini0893.ElectrostaticEnvironment}
  Electrostatic environment used by the ideal point-charge model
\end{definition}
\begin{proof}
  This is the declared model definition of Electrostatic Environment.
\end{proof}
\begin{definition}[Coordinate Convention]
  \label{def:physics:phyx-mini-0893:phyxminiproblems-problemphyxmini0893-coordinateconvention}
  \lean{PhyXMiniProblems.ProblemPhyXMini0893.CoordinateConvention}
  Coordinate convention for components of the dimensionless directions
\end{definition}
\begin{proof}
  This is the declared model definition of Coordinate Convention.
\end{proof}
\begin{definition}[Three Charge Triangle Figure]
  \label{def:physics:phyx-mini-0893:phyxminiproblems-problemphyxmini0893-threechargetrianglefigure}
  \lean{PhyXMiniProblems.ProblemPhyXMini0893.ThreeChargeTriangleFigure}
  \uses{def:physics:phyx-mini-0893:phyxminiproblems-problemphyxmini0893-trianglevertex, def:physics:phyx-mini-0893:phyxminiproblems-problemphyxmini0893-triangleedge}
  Literal presentation data from 893.png. Numerical fields are explicitly unit-labelled scalar readouts from the raster and do not replace the physical quantities in EquilateralThreeChargeSetup
\end{definition}
\begin{proof}
  This is the declared model definition of Three Charge Triangle Figure.
\end{proof}
\begin{definition}[Equilateral Three Charge Setup]
  \label{def:physics:phyx-mini-0893:phyxminiproblems-problemphyxmini0893-equilateralthreechargesetup}
  \lean{PhyXMiniProblems.ProblemPhyXMini0893.EquilateralThreeChargeSetup}
  \uses{def:physics:phyx-mini-0893:phyxminiproblems-problemphyxmini0893-planarvector, def:physics:phyx-mini-0893:phyxminiproblems-problemphyxmini0893-lengthquantity, def:physics:phyx-mini-0893:phyxminiproblems-problemphyxmini0893-chargemagnitudequantity, def:physics:phyx-mini-0893:phyxminiproblems-problemphyxmini0893-forcemagnitudequantity, def:physics:phyx-mini-0893:phyxminiproblems-problemphyxmini0893-forcevectorquantity, def:physics:phyx-mini-0893:phyxminiproblems-problemphyxmini0893-coulombconstantquantity, def:physics:phyx-mini-0893:phyxminiproblems-problemphyxmini0893-trianglevertex, def:physics:phyx-mini-0893:phyxminiproblems-problemphyxmini0893-triangleedge, def:physics:phyx-mini-0893:phyxminiproblems-problemphyxmini0893-lowercharge, def:physics:phyx-mini-0893:phyxminiproblems-problemphyxmini0893-chargesign, def:physics:phyx-mini-0893:phyxminiproblems-problemphyxmini0893-electrostaticenvironment, def:physics:phyx-mini-0893:phyxminiproblems-problemphyxmini0893-coordinateconvention, def:physics:phyx-mini-0893:phyxminiproblems-problemphyxmini0893-threechargetrianglefigure}
  Independent physical data for the configuration. The individual force vectors, net force vector, and net-force magnitude are fields constrained by the laws below; no field is assigned a displayed answer value
\end{definition}
\begin{proof}
  This is the declared model definition of Equilateral Three Charge Setup.
\end{proof}
\begin{definition}[Matches Three Positive Point Charge Scenario]
  \label{def:physics:phyx-mini-0893:phyxminiproblems-problemphyxmini0893-matchesthreepositivepointchargescenario}
  \lean{PhyXMiniProblems.ProblemPhyXMini0893.MatchesThreePositivePointChargeScenario}
  \uses{def:physics:phyx-mini-0893:phyxminiproblems-problemphyxmini0893-equilateralthreechargesetup}
  Qualitative point-charge and sign assumptions stated by the problem
\end{definition}
\begin{proof}
  This is the declared model definition of Matches Three Positive Point Charge Scenario.
\end{proof}
\begin{definition}[Matches Supplied Three Charge Triangle Figure]
  \label{def:physics:phyx-mini-0893:phyxminiproblems-problemphyxmini0893-matchessuppliedthreechargetrianglefigure}
  \lean{PhyXMiniProblems.ProblemPhyXMini0893.MatchesSuppliedThreeChargeTriangleFigure}
  \uses{def:physics:phyx-mini-0893:phyxminiproblems-problemphyxmini0893-lengthincentimeters, def:physics:phyx-mini-0893:phyxminiproblems-problemphyxmini0893-chargemagnitudeinnanocoulombs, def:physics:phyx-mini-0893:phyxminiproblems-problemphyxmini0893-equilateralthreechargesetup}
  All label, glyph, and geometry evidence read from the primary raster. The top angle is not explicitly drawn, so its raster fields are false and none; its physical 60 degree value follows from equilateral geometry
\end{definition}
\begin{proof}
  This is the declared model definition of Matches Supplied Three Charge Triangle Figure.
\end{proof}
\begin{definition}[Uses Equilateral Triangle Geometry]
  \label{def:physics:phyx-mini-0893:phyxminiproblems-problemphyxmini0893-usesequilateraltrianglegeometry}
  \lean{PhyXMiniProblems.ProblemPhyXMini0893.UsesEquilateralTriangleGeometry}
  \uses{def:physics:phyx-mini-0893:phyxminiproblems-problemphyxmini0893-lengthinmeters, def:physics:phyx-mini-0893:phyxminiproblems-problemphyxmini0893-equilateralthreechargesetup}
  The equilateral geometry and the two unit directions from the lower vertices toward the top vertex. Component zero points right and component one points up. These are figure-derived geometric relations, not force magnitudes
\end{definition}
\begin{proof}
  This is the declared model definition of Uses Equilateral Triangle Geometry.
\end{proof}
\begin{definition}[Has Physical Three Charge Parameters]
  \label{def:physics:phyx-mini-0893:phyxminiproblems-problemphyxmini0893-hasphysicalthreechargeparameters}
  \lean{PhyXMiniProblems.ProblemPhyXMini0893.HasPhysicalThreeChargeParameters}
  \uses{def:physics:phyx-mini-0893:phyxminiproblems-problemphyxmini0893-lengthinmeters, def:physics:phyx-mini-0893:phyxminiproblems-problemphyxmini0893-chargemagnitudeincoulombs, def:physics:phyx-mini-0893:phyxminiproblems-problemphyxmini0893-coulombconstantinnewtonmeterssquaredpercoulombsquared, def:physics:phyx-mini-0893:phyxminiproblems-problemphyxmini0893-equilateralthreechargesetup}
  Positivity and nondegeneracy conditions for the physical parameters
\end{definition}
\begin{proof}
  This is the declared model definition of Has Physical Three Charge Parameters.
\end{proof}
\begin{definition}[Uses Textbook Vacuum Coulomb Constant]
  \label{def:physics:phyx-mini-0893:phyxminiproblems-problemphyxmini0893-usestextbookvacuumcoulombconstant}
  \lean{PhyXMiniProblems.ProblemPhyXMini0893.UsesTextbookVacuumCoulombConstant}
  \uses{def:physics:phyx-mini-0893:phyxminiproblems-problemphyxmini0893-coulombconstantinnewtonmeterssquaredpercoulombsquared, def:physics:phyx-mini-0893:phyxminiproblems-problemphyxmini0893-equilateralthreechargesetup}
  The unit-independent Coulomb constant is connected to Physlib's Electromagnetism.EMSystem.coulombConstant. The rounded value is the standard textbook calibration and contains no requested force magnitude
\end{definition}
\begin{proof}
  This is the declared model definition of Uses Textbook Vacuum Coulomb Constant.
\end{proof}
\begin{definition}[Satisfies Coulomb Force Law And Superposition]
  \label{def:physics:phyx-mini-0893:phyxminiproblems-problemphyxmini0893-satisfiescoulombforcelawandsuperposition}
  \lean{PhyXMiniProblems.ProblemPhyXMini0893.SatisfiesCoulombForceLawAndSuperposition}
  \uses{def:physics:phyx-mini-0893:phyxminiproblems-problemphyxmini0893-lengthreadout, def:physics:phyx-mini-0893:phyxminiproblems-problemphyxmini0893-chargemagnitudereadout, def:physics:phyx-mini-0893:phyxminiproblems-problemphyxmini0893-forcemagnitudereadout, def:physics:phyx-mini-0893:phyxminiproblems-problemphyxmini0893-forcevectorreadout, def:physics:phyx-mini-0893:phyxminiproblems-problemphyxmini0893-coulombconstantreadout, def:physics:phyx-mini-0893:phyxminiproblems-problemphyxmini0893-lowercharge, def:physics:phyx-mini-0893:phyxminiproblems-problemphyxmini0893-equilateralthreechargesetup}
  Coulomb's repulsive point-charge law in vector form, force superposition, and the defining norm relation between the net vector and its magnitude. These equations contain no numerical value of the requested net force
\end{definition}
\begin{proof}
  This is the declared model definition of Satisfies Coulomb Force Law And Superposition.
\end{proof}
\begin{definition}[net Force Magnitude In Newtons]
  \label{def:physics:phyx-mini-0893:phyxminiproblems-problemphyxmini0893-netforcemagnitudeinnewtons}
  \lean{PhyXMiniProblems.ProblemPhyXMini0893.netForceMagnitudeInNewtons}
  \uses{def:physics:phyx-mini-0893:phyxminiproblems-problemphyxmini0893-forcemagnitudeinnewtons, def:physics:phyx-mini-0893:phyxminiproblems-problemphyxmini0893-equilateralthreechargesetup}
  The requested magnitude of the net force on the top charge, in newtons
\end{definition}
\begin{proof}
  This is the declared model definition of net Force Magnitude In Newtons.
\end{proof}
\begin{definition}[Answer Choice]
  \label{def:physics:phyx-mini-0893:phyxminiproblems-problemphyxmini0893-answerchoice}
  \lean{PhyXMiniProblems.ProblemPhyXMini0893.AnswerChoice}
  The four answer labels printed with the problem
\end{definition}
\begin{proof}
  This is the declared model definition of Answer Choice.
\end{proof}
\begin{definition}[displayed Force In Newtons]
  \label{def:physics:phyx-mini-0893:phyxminiproblems-problemphyxmini0893-answerchoice-displayedforceinnewtons}
  \lean{PhyXMiniProblems.ProblemPhyXMini0893.AnswerChoice.displayedForceInNewtons}
  \uses{def:physics:phyx-mini-0893:phyxminiproblems-problemphyxmini0893-answerchoice}
  The force magnitude printed beside an answer choice, in newtons
\end{definition}
\begin{proof}
  This is the declared model definition of displayed Force In Newtons.
\end{proof}
\begin{definition}[recorded Dataset Answer]
  \label{def:physics:phyx-mini-0893:phyxminiproblems-problemphyxmini0893-recordeddatasetanswer}
  \lean{PhyXMiniProblems.ProblemPhyXMini0893.recordedDatasetAnswer}
  \uses{def:physics:phyx-mini-0893:phyxminiproblems-problemphyxmini0893-answerchoice}
  The answer label recorded by the source dataset
\end{definition}
\begin{proof}
  This is the declared model definition of recorded Dataset Answer.
\end{proof}
\begin{definition}[Rounds To Nearest Resolution]
  \label{def:physics:phyx-mini-0893:phyxminiproblems-problemphyxmini0893-roundstonearestresolution}
  \lean{PhyXMiniProblems.ProblemPhyXMini0893.RoundsToNearestResolution}
  A physical value rounds to displayed at the stated positive resolution
\end{definition}
\begin{proof}
  This is the declared model definition of Rounds To Nearest Resolution.
\end{proof}
\begin{definition}[Is Unique Closest Displayed Answer]
  \label{def:physics:phyx-mini-0893:phyxminiproblems-problemphyxmini0893-isuniqueclosestdisplayedanswer}
  \lean{PhyXMiniProblems.ProblemPhyXMini0893.IsUniqueClosestDisplayedAnswer}
  \uses{def:physics:phyx-mini-0893:phyxminiproblems-problemphyxmini0893-equilateralthreechargesetup, def:physics:phyx-mini-0893:phyxminiproblems-problemphyxmini0893-netforcemagnitudeinnewtons, def:physics:phyx-mini-0893:phyxminiproblems-problemphyxmini0893-answerchoice}
  A choice is uniquely closest to the calculated physical force magnitude
\end{definition}
\begin{proof}
  This is the declared model definition of Is Unique Closest Displayed Answer.
\end{proof}
\begin{lemma}[net Force Magnitude rounding Bounds]
  \label{lem:physics:phyx-mini-0893:phyxminiproblems-problemphyxmini0893-netforcemagnitude-roundingbounds}
  \lean{PhyXMiniProblems.ProblemPhyXMini0893.netForceMagnitude_roundingBounds}
  \uses{def:physics:phyx-mini-0893:phyxminiproblems-problemphyxmini0893-equilateralthreechargesetup, def:physics:phyx-mini-0893:phyxminiproblems-problemphyxmini0893-matchesthreepositivepointchargescenario, def:physics:phyx-mini-0893:phyxminiproblems-problemphyxmini0893-matchessuppliedthreechargetrianglefigure, def:physics:phyx-mini-0893:phyxminiproblems-problemphyxmini0893-usesequilateraltrianglegeometry, def:physics:phyx-mini-0893:phyxminiproblems-problemphyxmini0893-hasphysicalthreechargeparameters, def:physics:phyx-mini-0893:phyxminiproblems-problemphyxmini0893-usestextbookvacuumcoulombconstant, def:physics:phyx-mini-0893:phyxminiproblems-problemphyxmini0893-satisfiescoulombforcelawandsuperposition, def:physics:phyx-mini-0893:phyxminiproblems-problemphyxmini0893-netforcemagnitudeinnewtons}
  The two equal pairwise forces have opposing horizontal components and equal upward components. Coulomb's law and the equilateral geometry therefore put the net magnitude between 3.05 × 10⁻⁴ N and 3.15 × 10⁻⁴ N
\end{lemma}
\begin{proof}
  The conclusion follows from the governing relations and figure readouts named in the statement of net Force Magnitude rounding Bounds.
\end{proof}
% --- Archon declaration topology end ---
% --- Archon physics formalization source end ---
